\documentclass[11pt,a4paper]{article}
\usepackage[top=2.3cm,bottom=2.3cm,left=2.2cm,right=2.2cm]{geometry}

\usepackage{amsmath,amssymb}
\usepackage{fontspec}
\setmainfont{Liberation Serif}
\setsansfont{Liberation Sans}
\setmonofont{DejaVu Sans Mono}[Scale=0.84]

\usepackage{xcolor}
\definecolor{brandblue}{RGB}{20, 55, 105}
\definecolor{accentblue}{RGB}{35, 95, 165}
\definecolor{codebg}{RGB}{248, 249, 251}
\definecolor{codeframe}{RGB}{215, 222, 230}
\definecolor{javakeyword}{RGB}{0, 45, 170}
\definecolor{javastring}{RGB}{5, 120, 20}
\definecolor{javacomment}{RGB}{120, 120, 120}
\definecolor{javanumber}{RGB}{20, 75, 220}

\definecolor{noteborder}{RGB}{30, 110, 65}
\definecolor{notebg}{RGB}{242, 249, 245}
\definecolor{warnborder}{RGB}{185, 75, 10}
\definecolor{warnbg}{RGB}{254, 248, 240}
\definecolor{tipborder}{RGB}{30, 95, 160}
\definecolor{tipbg}{RGB}{242, 247, 254}

\usepackage{fancyhdr}
\usepackage{lastpage}
\pagestyle{fancy}
\fancyhf{}
\lhead{\parbox[t]{0.58\textwidth}{\raggedright\small\sffamily\color{brandblue}\textbf{T01 --- Course Introduction}}}
\rhead{\small\sffamily\color{gray}Programmazione II -- UniVR (2025/2026)}
\cfoot{\small\sffamily Pagina \thepage\ di \pageref{LastPage}}
\renewcommand{\headrulewidth}{0.5pt}
\renewcommand{\headrule}{\hbox to\headwidth{\color{brandblue!70!black}\leaders\hrule height \headrulewidth\hfill}}

\usepackage{titlesec}
\titleformat{\section}{\Large\bfseries\sffamily\color{brandblue}}{\thesection}{0.8em}{}[\titlerule]
\titleformat{\subsection}{\large\bfseries\sffamily\color{accentblue}}{\thesubsection}{0.8em}{}
\titleformat{\subsubsection}{\normalsize\bfseries\sffamily\color{brandblue!85!black}}{\thesubsubsection}{0.8em}{}

\usepackage{needspace}
\usepackage{fancyvrb}
\DefineVerbatimEnvironment{asciibox}{Verbatim}{
  fontsize=\fontsize{7.5pt}{9.2pt}\selectfont,
  frame=single,
  rulecolor=\color{codeframe},
  fillcolor=\color{codebg},
  framesep=2.5mm
}
\usepackage{listings}
\lstset{
  backgroundcolor=\color{codebg},
  frame=single,
  rulecolor=\color{codeframe},
  basicstyle=\ttfamily\footnotesize,
  keywordstyle=\color{javakeyword}\bfseries,
  stringstyle=\color{javastring},
  commentstyle=\color{javacomment}\itshape,
  numberstyle=\tiny\color{gray},
  numbers=left,
  numbersep=7pt,
  stepnumber=1,
  breaklines=true,
  breakatwhitespace=true,
  showstringspaces=false,
  tabsize=4,
  xleftmargin=12pt,
  framexleftmargin=12pt
}

\newcommand{\passthrough}[1]{#1}
\providecommand{\tightlist}{\setlength{\itemsep}{0pt}\setlength{\parskip}{0pt}}

\usepackage{tcolorbox}
\tcbuselibrary{skins,breakable}
\newtcolorbox{studynote}[1][Nota]{
  enhanced,
  breakable,
  colback=notebg,
  colframe=noteborder,
  fonttitle=\bfseries\sffamily\small,
  coltitle=white,
  title={#1},
  arc=2mm,
  boxrule=0.8pt,
  left=9pt, right=9pt, top=7pt, bottom=7pt
}

\newtcolorbox{studywarning}[1][Attenzione]{
  enhanced,
  breakable,
  colback=warnbg,
  colframe=warnborder,
  fonttitle=\bfseries\sffamily\small,
  coltitle=white,
  title={#1},
  arc=2mm,
  boxrule=0.8pt,
  left=9pt, right=9pt, top=7pt, bottom=7pt
}

\newtcolorbox{studytip}[1][Suggerimento]{
  enhanced,
  breakable,
  colback=tipbg,
  colframe=tipborder,
  fonttitle=\bfseries\sffamily\small,
  coltitle=white,
  title={#1},
  arc=2mm,
  boxrule=0.8pt,
  left=9pt, right=9pt, top=7pt, bottom=7pt
}

\usepackage{booktabs}
\usepackage{longtable}
\usepackage{array}
\usepackage{calc}

\usepackage{hyperref}
\hypersetup{
  colorlinks=true,
  linkcolor=accentblue,
  urlcolor=accentblue,
  citecolor=brandblue,
  pdfborder={0 0 0}
}

\title{\Huge\bfseries\sffamily\color{brandblue} T01 --- Course Introduction \& Metodologia del Corso \\[0.3em] \large\sffamily Obiettivi Formativi, Modalità d'Esame al Calcolatore, Toolchain e Syllabus}
\author{\textbf{Docente:} Prof. Michele Pasqua \quad---\quad \textbf{Appunti Revisione:} Wally}
\date{A.A. 2025/2026 -- Universit\`a degli Studi di Verona}

\begin{document}
\maketitle
\thispagestyle{fancy}
\vspace{-1.2em}
\noindent\textcolor{brandblue}{\rule{\textwidth}{0.8pt}}
\vspace{1.2em}

Benvenuto in questa revisione completa e metodica di
\textbf{Programmazione II} (Prof.~Michele Pasqua, UniVR 2025/2026).
Seguiremo fedelmente le regole stabilite: \textbf{nessuna omissione},
analisi minuziosa di ogni concetto teorico, confronto riga per riga del
codice, evidenziazione delle finezze d'esame e stop interattivo a ogni
blocco.

\begin{center}\rule{0.5\linewidth}{0.5pt}\end{center}

\subsubsection{\texorpdfstring{1. Teoria: Concetti Teorici dalle Slide
(\href{file:///home/wally/Documenti/GitHub/Programmazione-II/25-26/Teoria-e-Esercitazioni/Slides-20260902/T01\%20-\%20Course\%20Introduction.pdf}{T01
- Course
Introduction.pdf})}{1. Teoria: Concetti Teorici dalle Slide (T01 - Course Introduction.pdf)}}\label{teoria-concetti-teorici-dalle-slide-t01---course-introduction.pdf}

La slide introduttiva definisce il perimetro accademico, gli obiettivi
formativi e le modalità d'esame:

\begin{enumerate}
\def\labelenumi{\arabic{enumi}.}
\tightlist
\item
  \textbf{Docente e Background Scientifico}:

  \begin{itemize}
  \tightlist
  \item
    \textbf{Prof.~Michele Pasqua} (Dipartimento di Informatica, Cà
    Vignal 2, 2° piano, Stanza 24 ---
    \passthrough{\lstinline!michele.pasqua@univr.it!}).
  \item
    Aree di ricerca: \emph{Program verification and validation},
    \emph{Programming language design and implementation}, \emph{Code
    protection and software security}.
  \item
    Laboratori e progetti: co-fondatore del laboratorio
    \textbf{UNIVERSE} (\emph{University of Verona Software Engineering
    Lab}, diretto dal Prof.~Mariano Ceccato) e co-fondatore dello
    spin-off universitario \textbf{KonTrust Srl} (tracciamento IoT su
    blockchain e infrastrutture anti-manomissione); referente per Verona
    del programma nazionale di formazione e gare su cybersecurity
    (\emph{CyberChallenge.IT}).
  \end{itemize}
\item
  \textbf{Organizzazione della Didattica}:

  \begin{itemize}
  \tightlist
  \item
    \textbf{Modulo di Teoria (40+ ore)}: Lunedì (16:30) e Martedì
    (13:30) in Aula A (2 ore ciascuna).
  \item
    \textbf{Modulo di Laboratorio (12+ ore)}: Martedì (10:30) in Aula
    Delta (2 ore).
  \item
    \textbf{Materiale didattico di riferimento}:

    \begin{itemize}
    \tightlist
    \item
      Paul J. Deitel \& Harvey M. Deitel, \emph{Java -- How to Program},
      Pearson (11/e -- 2018).
    \item
      Claudio De Sio Cesari, \emph{Il nuovo Java -- Guida alla
      programmazione moderna}, Hoepli (1/e -- 2020).
    \item
      Moodle UniVR: slide, esercizi di laboratorio e soluzioni,
      simulazioni.
    \end{itemize}
  \end{itemize}
\item
  \textbf{Modalità d'Esame e Valutazione}:

  \begin{itemize}
  \tightlist
  \item
    \textbf{Formato}: Prova pratica al calcolatore in \textbf{Aula
    Delta} (ambiente Linux/IDE).
  \item
    \textbf{Contenuto}: Risoluzione di una serie di esercizi di
    programmazione in linguaggio Java (sviluppo classi, gerarchie a
    oggetti, implementazione interfacce, algoritmi su collezioni,
    gestione eccezioni, test con asserzioni/JUnit).
  \item
    \textbf{Voto finale}: Somma algebrica dei punti ottenuti nei singoli
    esercizi.
  \item
    È prevista una \textbf{simulazione d'esame} ufficiale a fine corso.
  \end{itemize}
\item
  \textbf{Syllabus del Corso (Mappa degli argomenti da T01 a T20)}:

  \begin{itemize}
  \tightlist
  \item
    \emph{Object-Orientation Principles} (T02)
  \item
    \emph{Building, Running, and Deployment} (T03)
  \item
    \emph{Types, Classes, and Objects} (T04)
  \item
    \emph{Java Syntax and Semantics} (T05)
  \item
    \emph{Java Classes and Objects} (T06)
  \item
    \emph{Libraries and String Class} (T07)
  \item
    \emph{Class Constructors and Encapsulation} (T08)
  \item
    \emph{Arrays and Enumerative Types} (T09)
  \item
    \emph{Java Variables and Call Stack} (T10)
  \item
    \emph{Class and Package Visibility} (T11)
  \item
    \emph{Inheritance and Polymorphism} (T12)
  \item
    \emph{Abstract Classes and Interfaces} (T13)
  \item
    \emph{Nested and Anonymous Classes} (T14)
  \item
    \emph{Generic Types} (T15)
  \item
    \emph{Collections} (T16)
  \item
    \emph{Error Handling with Exceptions} (T17)
  \item
    \emph{Functional Interfaces and Streams} (T18)
  \item
    \emph{Documentation and Unit Testing} (T19)
  \item
    \emph{Java IO} (T20)
  \end{itemize}
\end{enumerate}

\begin{center}\rule{0.5\linewidth}{0.5pt}\end{center}

\subsubsection{2. Codice: Analisi del Codice della Lezione
1}\label{codice-analisi-del-codice-della-lezione-1}

\begin{itemize}
\tightlist
\item
  Nella cartella ufficiale delle lezioni
  \href{file:///home/wally/Documenti/GitHub/Programmazione-II/25-26/Teoria-e-Esercitazioni/Codice-20260902}{\passthrough{\lstinline!Codice-20260902!}}:

  \begin{itemize}
  \tightlist
  \item
    \textbf{Non sono presenti file di codice per le Lezioni 01, 02 e
    03}.
  \item
    Il primo codice sorgente ufficiale del docente compare a partire
    dalla \textbf{Lezione 04} con
    \href{file:///home/wally/Documenti/GitHub/Programmazione-II/25-26/Teoria-e-Esercitazioni/Codice-20260902/Lezione04/Example.java}{\passthrough{\lstinline!Example.java!}}
    (dedicato a tipi primitivi, reference, casting e operatori), mentre
    il progetto a oggetti incrementale vero e proprio parte dal blocco
    \textbf{Lezioni 05-06} con il progetto
    \passthrough{\lstinline!SimpleDate!}
    (\passthrough{\lstinline!Date.java!} e
    \passthrough{\lstinline!MainDate.java!}).
  \end{itemize}
\item
  Nel tuo workspace
  \href{file:///home/wally/Documenti/GitHub/Programmazione-II/25-26/esercizi-wally-25-26}{\passthrough{\lstinline!esercizi-wally-25-26!}}:

  \begin{itemize}
  \tightlist
  \item
    La cartella
    \href{file:///home/wally/Documenti/GitHub/Programmazione-II/25-26/esercizi-wally-25-26/es01}{\passthrough{\lstinline!es01!}}
    è già predisposta come modulo Maven con la prima incarnazione di
    \passthrough{\lstinline!Date!} e \passthrough{\lstinline!MainDate!},
    sincronizzata con la struttura multi-modulo orchestrata dal
    \href{file:///home/wally/Documenti/GitHub/Programmazione-II/25-26/esercizi-wally-25-26/pom.xml}{\passthrough{\lstinline!pom.xml!}}
    di root e dallo script
    \href{file:///home/wally/Documenti/GitHub/Programmazione-II/25-26/esercizi-wally-25-26/esercizi.sh}{\passthrough{\lstinline!esercizi.sh!}}.
  \end{itemize}
\end{itemize}

\begin{center}\rule{0.5\linewidth}{0.5pt}\end{center}

\subsubsection{3. Setup: ️ Configurazione Operativa
dell'Ambiente}\label{setup-configurazione-operativa-dellambiente}

Per garantire che tutti i passaggi da qui in avanti siano riproducibili
su IntelliJ IDEA e da CLI: 1. Verifica che la JDK attiva sia Java 17 o
superiore: \passthrough{\lstinline!bash    java -version    mvn -v!} 2.
Il file di root
\href{file:///home/wally/Documenti/GitHub/Programmazione-II/25-26/esercizi-wally-25-26/pom.xml}{\passthrough{\lstinline!pom.xml!}}
gestisce l'aggregazione di tutti i sottomoduli. Per ogni blocco
successivo, quando andremo ad analizzare le classi introdotte a lezione,
verificheremo che compilino ed eseguano pulitamente sia tramite:
\passthrough{\lstinline!bash    ./esercizi.sh run <numero\_modulo>!} sia
all'interno di IntelliJ IDEA.

\begin{center}\rule{0.5\linewidth}{0.5pt}\end{center}

\begin{studynote}[Nota di Studio] La \textbf{Lezione 1 / T01} esaurisce qui la parte
introduttiva/organizzativa. Il prossimo blocco è la \textbf{Lezione 2 /
Slide T02: \emph{Object-Orientation and UML}}, dove entreremo nel vivo
con: - Storia e tassonomia dei paradigmi: \emph{Imperativo/Procedurale},
\emph{Dichiarativo/Funzionale}, \emph{Object-Oriented}. - I 4 pilastri
dell'OOP: \emph{Astrazione, Incapsulamento, Ereditarietà, Polimorfismo}.
- Concetto di \emph{Classe} vs \emph{Oggetto} (stato, comportamento,
identità). - Notazione UML per classi e associazioni (\emph{Visibility
markers}, cardinalità, tipi di relazione).
\end{studynote}

\textbf{Dammi conferma quando sei pronto per passare a T02 (procederemo
analizzando il blocco concettuale un po' alla volta)!}


\end{document}
